%!TEX program = lualatex
\documentclass{loopmind-doc}

\doclang{es}
\doctype{Documento de producto}
\doctitle{Makis\newline De URL a workspace}
\docsubtitle{El onboarding autónomo que convierte el sitio de una empresa en un sistema de trabajo de GTM.}
\docclient{Equipo Maki Acevichado}
\docdate{Septiembre 2026}
\docversion{1.0}

\usetikzlibrary{positioning,arrows.meta,shapes.geometric}
\setlength{\headheight}{20pt}
\newcommand{\makisstep}[3]{%
\begin{tcolorbox}[colback=lm-surface,colframe=lm-border,boxrule=.45pt,arc=2mm,
left=4mm,right=4mm,top=4mm,bottom=4mm]
{\datafont\bfseries\fontsize{18pt}{18pt}\selectfont\textcolor{lm-accent}{#1}}\hspace{7pt}%
{\headingfont\bfseries\fontsize{11pt}{13pt}\selectfont\color{lm-black}#2}\par\vspace{4pt}
{\fontsize{9pt}{13pt}\selectfont\color{lm-secondary}#3}
\end{tcolorbox}}

\begin{document}

\makecover
\sectionlabel{La idea}

\section{Una puerta de entrada, no un formulario}

Makis empieza con la evidencia que ya tiene la empresa: su sitio web. El usuario pega una URL, confirma lo que Makis entendió y recibe un workspace de Go-To-Market listo para revisar.

\begin{callouttitled}{La decisión de producto}
La documentación vive en Markdown porque debe ser durable, versionable y exportable. La interfaz la renderiza como vistas claras para que el usuario pueda decidir y actuar sin leer archivos técnicos.
\end{callouttitled}

\vspace{8pt}
\begin{center}
\begin{tikzpicture}[node distance=4mm and 5mm, every node/.style={align=center}]
\node[rounded corners=3pt, fill=lm-accent, text=white, minimum width=28mm, minimum height=16mm, font=\bfseries\small] (url) {Pega\\su URL};
\node[rounded corners=3pt, fill=lm-accent-tint, draw=lm-accent, minimum width=31mm, minimum height=16mm, font=\bfseries\small, right=of url] (read) {Makis lee\\el sitio};
\node[rounded corners=3pt, fill=lm-surface, draw=lm-border, minimum width=31mm, minimum height=16mm, font=\bfseries\small, right=of read] (confirm) {Humano\\confirma};
\node[rounded corners=3pt, fill=lm-black, text=white, minimum width=35mm, minimum height=16mm, font=\bfseries\small, right=of confirm] (work) {Workspace\\listo};
\draw[-{Latex[length=2mm]}, thick, color=lm-accent] (url) -- (read);
\draw[-{Latex[length=2mm]}, thick, color=lm-accent] (read) -- (confirm);
\draw[-{Latex[length=2mm]}, thick, color=lm-accent] (confirm) -- (work);
\end{tikzpicture}
\end{center}

\section{El recorrido completo}

\makisstep{01}{Descubrir}{La URL, el objetivo y el presupuesto mínimo crean un brief verificable. Makis extrae marca, oferta, tono, categorías y llamados a la acción.}
\makisstep{02}{Investigar}{Exa y la lectura del sitio ubican competidores, tendencias, lenguaje de categoría y oportunidades.}
\makisstep{03}{Construir}{Makis crea la estrategia, el plan de contenido, las creatividades y los documentos que explican cada decisión.}
\makisstep{04}{Gobernar}{El usuario revisa en la UI, edita en línea, aprueba o pide una nueva ejecución. Nada importante se ejecuta sin control humano.}
\makisstep{05}{Aprender}{Las métricas y el resultado de cada pieza se convierten en un aprendizaje visible para el siguiente experimento.}

\section{El workspace que aparece solo}

\begin{center}
\begin{tikzpicture}[
folder/.style={rounded corners=2pt, fill=lm-surface, draw=lm-border, minimum width=114mm, minimum height=9mm, align=left, font=\small},
doc/.style={rounded corners=2pt, fill=white, draw=lm-border-light, minimum width=102mm, minimum height=8mm, align=left, font=\small}
]
\node[folder] (root) {\textbf{workspace / empresa} \hfill contexto y decisiones del negocio};
\node[doc, below=3mm of root] (brief) {00-brief.md \hfill perfil confirmado, objetivo y supuestos};
\node[doc, below=2mm of brief] (research) {01-investigacion.md \hfill fuentes, competidores y hallazgos};
\node[doc, below=2mm of research] (strategy) {02-estrategia.md \hfill audiencia, oferta, canales y KPIs};
\node[doc, below=2mm of strategy] (plan) {03-plan-de-contenido.md \hfill piezas, prioridades y calendario};
\node[doc, below=2mm of plan] (creative) {04-creatividades.md \hfill copy, prompts y aprobaciones};
\node[doc, below=2mm of creative] (results) {05-resultados.md \hfill métricas, aprendizaje y siguiente experimento};
\end{tikzpicture}
\end{center}

Cada documento tiene fecha, fuentes, versión y estado de aprobación. SQLite no reemplaza estos archivos: indexa versiones, decisiones y eventos para que el sistema pueda recuperarlos y aprender.

\section{Una fuente de verdad, dos formas de verla}

\begin{center}
\begin{tikzpicture}[node distance=10mm, every node/.style={align=center}]
\node[rounded corners=3pt, fill=lm-black, text=white, minimum width=44mm, minimum height=22mm, font=\bfseries\small] (md) {Markdown\\fuente de verdad};
\node[rounded corners=3pt, fill=lm-accent, text=white, minimum width=44mm, minimum height=22mm, font=\bfseries\small, right=of md] (ui) {UI renderizada\\para decidir};
\draw[-{Latex[length=2.5mm]}, thick, color=lm-accent] (md) -- node[above, font=\scriptsize\color{lm-secondary}]{renderiza} (ui);
\draw[-{Latex[length=2.5mm]}, thick, color=lm-accent] (ui.south) to[out=-100,in=-80,looseness=1.1] node[below, font=\scriptsize\color{lm-secondary}]{un cambio aprobado actualiza y versiona} (md.south);
\end{tikzpicture}
\end{center}

\vspace{6pt}
\begin{tabularx}{\textwidth}{>{\bfseries}p{31mm}X X}
\toprule
& \textbf{Documento} & \textbf{Vista en producto} \\
\midrule
Brief & Datos del negocio y supuestos & Ficha para corregir contexto \\
Investigación & Fuentes y competidores & Tarjetas de evidencia \\
Estrategia & Audiencia, oferta y KPIs & Tablero de dirección \\
Creatividades & Copy, prompts e imágenes & Studio para editar y aprobar \\
Resultados & Métricas y aprendizaje & Gráfico y próximo experimento \\
\bottomrule
\end{tabularx}

\section{Qué debe probar el demo}

\begin{callout}
\begin{itemize}
  \item Una URL crea un perfil de negocio y un workspace visible en menos de un minuto.
  \item Tres documentos Markdown se generan y aparecen como módulos de interfaz, no como texto plano.
  \item Una corrección hecha en la UI queda registrada en el documento fuente.
  \item Una creatividad se aprueba y pasa a una ejecución real por Resend o a una simulación explícita.
  \item Un aprendizaje del dashboard define el siguiente experimento.
\end{itemize}
\end{callout}

\section{La frase de cierre}

\begin{tcolorbox}[colback=lm-black,colframe=lm-black,arc=2mm,left=7mm,right=7mm,top=7mm,bottom=7mm]
{\headingfont\bfseries\fontsize{17pt}{23pt}\selectfont\color{white}
Makis no pide que describas tu empresa a un agente. Entra por tu URL, construye el sistema de trabajo y te deja gobernar cada decisión.}
\end{tcolorbox}

\end{document}
